\documentclass[12pt,a4paper]{article}
\usepackage[T2A,T1]{fontenc}
\usepackage[utf8]{inputenc}
\usepackage[english,russian]{babel}
\usepackage{XCharter}
\usepackage{amsmath,amssymb,mathtools}
\usepackage{array,longtable,booktabs}
\usepackage[margin=24mm]{geometry}
\usepackage[hidelinks]{hyperref}
\urlstyle{same}
\setlength{\parindent}{0pt}
\setlength{\parskip}{0.45em}
\setlength{\LTleft}{0pt}
\setlength{\LTright}{0pt}
\emergencystretch=2em
\setlength{\tabcolsep}{4pt}
\newcolumntype{P}[1]{>{\raggedright\arraybackslash}p{#1}}
\title{Проверка идей Брусенцова в Lean: возможность и целесообразность}
\author{}
\date{}
\begin{document}
\maketitle

\emph{Заключение по состоянию на 20.11.2026. Это оценка метода, а не результат уже завершённой формализации.}

\section{Вывод}

\textbf{Lean пригоден и полезен для проверки формальных следствий из точно выписанной модели Брусенцова.} Он может проверить таблицы операций, свойства шкал, выводимость силлогизмов, наличие контрпримеров в конечных моделях и подсчёт модусов при заранее зафиксированных определениях.

\textbf{Сначала следует сделать небольшой проверочный скрипт; Lean оправдан на следующем этапе.} Скрипт быстрее ответит на вопросы о четырёх ячейках и нескольких конечных комбинациях. Lean становится полезнее, когда нужно доказать свойства для всех допустимых шкал или областей, записать условия явно, проверять зависимости теорем от аксиом и сохранить воспроизводимый доказательный файл.

Lean не решит, верно ли Брусенцов истолковал Аристотеля и следует ли называть его отношение «содержательным». Он проверит формализацию и доказательства \textbf{при выбранных определениях и предпосылках}.

\section{Что именно можно формализовать}

В главе «Реанимация аристотелевой силлогистики» (2005, электронные с. 2–5) заданы три статуса классовых выражений: необходимое включение \(+1\), привходящий статус \(0\), необходимое исключение \(-1\). Там же описаны операции min/max и смена знака, шкалы из \(2^n\) троичных значений, профиль \(+-0+\) и получение Barbara посредством соединения выражений и устранения среднего термина.

Lean можно использовать для таких проверок:

\begin{enumerate}
\item \textbf{Таблицы значений.} Задать конечный тип «трит» со значениями «исключено», «привходяще», «включено»; определить порядок, min/max и инверсию; доказать свойства операций или исчерпывающе проверить таблицы.
\item \textbf{Шкалы.} Представить шкалу из \(2^n\) тритов как конечную последовательность и проверить, что операции над шкалами совпадают с определениями автора.
\item \textbf{Четыре сочетания терминов.} Задать область \(U\), признаки \(X,Y\), четыре комбинации и постулат сосуществования. Отдельно задать авторское значение каждой позиции и отдельно — учебный экзистенциальный перевод профиля в условия «в этой группе есть предмет / таких предметов нет».
\item \textbf{Силлогизмы.} Формализовать соединение посылок, исключение среднего термина и Barbara. Доказать, что правило сохраняет выбранный статус или корректно для каждой модели, удовлетворяющей заданным условиям.
\item \textbf{Контрмодели и подсчёт.} Перечислить все оценки на конечной области, искать пример, где посылки выполнены, а заключение нет, и пересчитать модусы для явно заданного критерия корректности.
\end{enumerate}
Здесь особенно важно не смешать два объекта. Троичность Брусенцова следует ввести отдельным типом модели, условно называемым \emph{Trit}. Тип \emph{Prop} в Lean будет языком утверждений \textbf{о} тритах, шкалах, классах и правилах. Тогда метаязык помощника не подменит исследуемую троичную модель классической таблицей истинности \(P\to Q\).

Lean позволяет явно подключать классические принципы, например закон исключённого третьего, когда они нужны в метатеории. Их следует записывать отдельно от троичного статуса выражения в модели: наличие \(P\lor\neg P\) в окружении доказательства не задаёт, какой статус Брусенцов приписывает классу или ячейке.

Также следует разнести кодировки из разных работ. В статье 2003 года функция \(A(x,y)\) использует \(1,0,\sigma\), где ноль — антиприсущий статус, а \(\sigma\) — привходящий. В главе 2005 года автор описывает \(\sigma\) как третье значение классового выражения, затем использует знаковую шкалу \(+1,-1,0\), где ноль уже означает привходящий статус. В профиле \(+-0+\) этот статус стоит в позиции \(x'y\). Перед доказательством нужно закрепить, какую именно версию и семантику мы формализуем.

\section{Как поставить проверку}

\subsection{Шаг 1. Зафиксировать текстовую интерпретацию}

Для каждого термина выписать:

\begin{itemize}
\item что именно обозначают \(x,y\) и их противоположности;
\item в каком порядке идут четыре клетки;
\item когда клетка необходима, исключена или привходяща;
\item что означает постулат сосуществования;
\item как устроены операции, связывание посылок и устранение среднего термина.
\end{itemize}
Каждое определение сопоставить с изображением страницы первичного PDF. Если в статье есть две кодировки, не объединять их в одну без доказанного соответствия.

\subsection{Шаг 2. Проверить конечную алгебру}

Сначала реализовать конечные триты и буквальные таблицы операций из статьи. Проверить все таблицы исчерпывающим вычислением и затем доказать выбранные алгебраические свойства. Убедиться, что значение \(0\) не спутано с пустой клеткой, ложным высказыванием или отсутствием объекта: это три разные вещи.

\subsection{Шаг 3. Вынести существование в явные условия}

Постулат сосуществования лучше задавать параметром теоремы или свойством модели, а не безусловной аксиомой всей формальной системы. Тогда можно отдельно увидеть:

\begin{itemize}
\item какие выводы требуют непустоты терминов;
\item какие работают без неё;
\item как меняется результат при пустых классах;
\item какие дополнительные существования скрыты в учебном переводе знаков.
\end{itemize}
Это делает зависимость результата видимой. В Lean можно просмотреть аксиомы, от которых зависит теорема, командой \emph{\#print axioms}.

\subsection{Шаг 4. Сверить вывод и пересчёт модусов}

Для каждого правила сформулировать теорему в выбранной семантике и проверить её. Числа 128, 190 и 192 следует пересчитывать только после того, как точно заданы:

\begin{itemize}
\item список форм и конфигураций;
\item условия существования терминов;
\item критерий «правильного» модуса;
\item учёт перестановок, фигур и эквивалентных форм.
\end{itemize}
Lean сможет подтвердить итоговый подсчёт для конкретного списка и определения. Он не объяснит сам, почему статьи Брусенцова сообщают разные числа; для этого потребуется текстологически восстановить, что именно считалось в каждой работе.

\section{Граница между формальной и исторической проверкой}

Lean проверяет доказательство утверждения, которое исследователь уже записал. Он не проверяет автоматически, соответствует ли формализация статье, греческому тексту Аристотеля или интуиции «содержательности». Это разные вопросы:

\begin{longtable}{P{0.290\textwidth}P{0.620\textwidth}}
\toprule
\textbf{Вопрос} & \textbf{Что может дать Lean} \\
\midrule
\endfirsthead
\toprule
\textbf{Вопрос} & \textbf{Что может дать Lean} \\
\midrule
\endhead
Следует ли Barbara из формальных правил модели? & Доказательство или контрпример для записанных правил. \\
Зависит ли вывод от постулата сосуществования? & Отдельные теоремы с этим условием и без него. \\
Верно ли авторское число модусов при данном определении? & Проверяемый перебор и воспроизводимый результат. \\
Так ли мыслил Аристотель? & Ничего напрямую; нужны текст, перевод и исторический анализ. \\
Является ли это отношение содержательным или предпочтительным? & Ничего нормативного; это философский и методологический вопрос. \\
Полностью ли формализация передаёт смысл статьи? & Сам Lean этого не решает; определения нужно сверять с первоисточником. \\
\bottomrule
\end{longtable}

\section{Целесообразность и ограничения}

\textbf{Рекомендация — пилот после аудита обозначений, а не немедленная полная формализация.} Для четырёх клеток и небольших таблиц обычный перебор проще и дешевле. Пилот на Lean оправдан, если его цель сформулирована как одна из следующих:

\begin{itemize}
\item проверить алгебраические законы троичной шкалы;
\item доказать корректность Barbara в реконструированной модели;
\item определить точную зависимость выводов от существования терминов;
\item воспроизводимо пересчитать конечный список модусов;
\item хранить модель, определения и доказательства в одном проверяемом проекте.
\end{itemize}
Для финальной проверочной версии стоит предпочесть обычные доказательства или \emph{decide +kernel}. В режиме \emph{+native}/\emph{native\_decide} результат вычисляется скомпилированным кодом и добавляет аксиому, поэтому увеличивается доверенная база. Это может ускорить большие вычисления, но для небольшого исследовательского пилота такой компромисс обычно не нужен. Следует также фиксировать все дополнительные аксиомы; нельзя объявлять спорное условие истинным только ради завершения доказательства.

\textbf{Версия и локальная готовность.} Официальная справка, закреплённая для этого заключения, относится к Lean 4.36.0. На этой рабочей машине команды \emph{lean}, \emph{lake} и \emph{elan} не найдены в \emph{PATH}, поэтому пробный файл здесь не запускался. Это заключение о технической пригодности и плане проверки, а не отчёт о готовой Lean-реализации.

\section{Библиография к заключению}

\subsection{Материалы Брусенцова}

\begin{itemize}
\item Н. П. Брусенцов. \href{https://ternarycomp.cs.msu.ru/Papers/Rea.pdf}{«Реанимация аристотелевой силлогистики» (2005), глава из \emph{Реставрации логики}, с. 140–145}. \href{../sources/pdfs/Брусенцов_Реанимация-аристотелевой-силлогистики_2005-глава-из-Реставрации-логики.pdf}{PDF в корпусе}.
\item Н. П. Брусенцов. \href{https://computer-museum.ru/histussr/trilog0.htm}{«Трехзначная интерпретация силлогистики Аристотеля» (2003)}. \href{../sources/pdfs/Брусенцов_Трехзначная-интерпретация-силлогистики-Аристотеля_2003_версия-музея.pdf}{Производная PDF-копия}.
\item Н. П. Брусенцов. \href{https://ternarycomp.cs.msu.ru/Papers/Soder_sled.pdf}{«Содержательное следование, сосуществование противоположностей и формальные импликации» (2006)}. \href{../sources/pdfs/Брусенцов_Содержательное-следование-и-формальные-импликации_2006.pdf}{PDF в корпусе}.
\item Н. П. Брусенцов. \href{https://ternarycomp.cs.msu.ru/Papers/Alrek.pdf}{«Алгебраическая реконструкция силлогистики» (2008)}. \href{../sources/pdfs/Брусенцов_Алгебраическая-реконструкция-силлогистики_2008.pdf}{PDF в корпусе}.
\end{itemize}
\subsection{Официальная документация Lean}

\begin{itemize}
\item Lean Language Reference, версия \href{https://lean-lang.org/doc/reference/4.36.0/}{4.36.0}: устройство языка и ядра проверки доказательств.
\item Lean Language Reference, \href{https://lean-lang.org/doc/reference/4.36.0/ValidatingProofs/}{Validating a Lean Proof}: что именно гарантирует проверенное доказательство и как проверять его зависимости.
\item Lean Language Reference, \href{https://lean-lang.org/doc/reference/4.36.0/Axioms/}{Axioms}: аксиомы, классические принципы и аудит зависимостей.
\item Lean Language Reference, \href{https://lean-lang.org/doc/reference/4.36.0/Tactic-Proofs/Tactic-Reference/}{Tactic Reference: decide}: вычисление конечных разрешимых утверждений и различие режимов \emph{+kernel} и \emph{+native}.
\item \href{https://lean-lang.org/theorem_proving_in_lean4/Propositions-and-Proofs/}{Theorem Proving in Lean 4: Propositions and Proofs}: отношение между утверждениями и доказательствами в Lean.
\item Mathlib documentation, \href{https://leanprover-community.github.io/mathlib4_docs/Mathlib/Data/Finset/Defs.html}{finite sets (Finset)}: инструменты для конечных областей и перечислений.
\end{itemize}
\section{Связанные документы проекта}

\begin{itemize}
\item \href{../docs/Саммари-идей-Брусенцова.tex}{Саммари идей Брусенцова}.
\item \href{Ревизия-корпуса-Брусенцова.tex}{Ревизия корпуса по первичным материалам}.
\item \href{Полный-указатель-трудов-Брусенцова.md}{Полный указатель найденных работ}.
\item \href{Карта-работ-Брусенцова.tex}{Карта корпуса и доступности}.
\end{itemize}
\end{document}
